\documentclass[11pt,a4paper]{article}

% ---- Packages ----
\usepackage[margin=2.5cm]{geometry}
\usepackage[utf8]{inputenc}
\usepackage[T1]{fontenc}
\usepackage{lmodern}
\usepackage{amsmath,amssymb,amsthm}
\usepackage{graphicx}
\usepackage[dvipsnames]{xcolor}
\usepackage{booktabs}
\usepackage{tabularx}
\usepackage{multirow}
\usepackage{enumitem}
\usepackage{caption}
\usepackage[numbers,sort&compress]{natbib}
\usepackage{lineno}
\usepackage{fancyhdr}
\usepackage[hidelinks,colorlinks=true,linkcolor=blue!60!black,citecolor=blue!60!black,urlcolor=blue!60!black]{hyperref}
\usepackage{microtype}
\usepackage{mdframed}

% ---- Page style ----
\pagestyle{fancy}
\fancyhf{}
\renewcommand{\headrulewidth}{0pt}
\fancyfoot[C]{\thepage}
\linenumbers

% ---- Paragraph style ----
\setlength{\parindent}{0pt}
\setlength{\parskip}{6pt}

% ---- Custom commands ----
\newcommand{\norm}[1]{\left\|#1\right\|}
\newcommand{\eps}{\varepsilon}

\begin{document}

\begin{center}
{\LARGE\bfseries Supplementary Information}\\[8pt]
{\large Designing Any Imaging System from Natural Language:\\
Agent-Constrained Composition over a Finite Primitive Basis}\\[12pt]
{\normalsize Chengshuai Yang}\\[6pt]
\today
\end{center}

\vspace{12pt}

% ========================================================================
\section{S1. Key Results from the Companion Paper}
% ========================================================================

The companion paper~\cite{paperII} establishes two foundational results that this paper builds upon:

\textbf{Finite Primitive Basis (FPB) Theorem.}
Any imaging forward model across 170 modalities and 5 carrier families (photon, X-ray, electron, acoustic, spin/particle) can be decomposed as a directed acyclic graph (DAG) over 11 canonical primitives: $P$ (propagation), $M$ (modulation), $\Pi$ (projection), $F$ (Fourier encoding), $C$ (convolution), $\Sigma$ (accumulation), $D$ (detection), $S$ (sampling), $W$ (dispersion), $R$ (scattering), and $\Lambda$ (nonlinear response).  The representation error $\eps_{\text{FPB}} < 0.01$ for all modalities.

\textbf{Triad Decomposition Theorem.}
Every reconstruction failure in computational imaging is caused by exactly one of three mechanisms:
\begin{itemize}[nosep]
\item \textbf{Gate 1 (Recoverability)}: Insufficient measurement information (under-sampling, wrong geometry).
\item \textbf{Gate 2 (Carrier Budget)}: Insufficient signal-to-noise ratio (weak source, noisy detector).
\item \textbf{Gate 3 (Operator Mismatch)}: Forward model does not match the true physics (calibration drift, unmodeled effects).
\end{itemize}
The total MSE decomposes as $\text{MSE} \leq \text{MSE}^{(G_1)} + \text{MSE}^{(G_2)} + \text{MSE}^{(G_3)}$, guaranteeing diagnostic completeness.

% ========================================================================
\section{S2. Novel System Design Specifications}
% ========================================================================

The 10 novel system designs compose FPB primitives into chains not found in prior literature.  Each is specified via the 8-field \texttt{spec.md} format.  The notation $\Phi_z$ denotes depth-parameterised convolution: the $C$ primitive applied with a depth-dependent point spread function $\Phi(z)$, generated by wave propagation through a random phase diffuser with defocus modulation.

\subsection{S2.1. 2D Systems}

\textbf{Lensless imaging} ($C \to D$, compression 1:1).
A thin random phase mask creates a caustic PSF with broad spatial-frequency support.  Reconstruction inverts the convolution via ADMM+TV.  PSNR: 43.7\,dB, SSIM: 0.984.

\begin{verbatim}
modality: lensless_imaging
carrier: photon
geometry: single_shot, 128x128
object: 128x128 2D image
forward_model: Convolve(C, psf=phase_mask) -> Detect(D)
noise: Poisson I_0=5000 + Gaussian sigma=3
target: PSNR >= 40 dB
system_elements:
  source: broadband LED
  optics: random phase mask (feature_scale=2.5)
  detector: CMOS 128x128
\end{verbatim}

\textbf{SIM} ($M \to C \to D$, super-resolution).
Structured illumination modulates the sample with sinusoidal patterns, shifting high-frequency content into the passband.  Three pattern orientations at three phases yield 9 measurements for 2$\times$ resolution enhancement.  PSNR: 27.7\,dB.

\subsection{S2.2. 3D Systems (8:1 compression)}

\textbf{3D Lensless} ($\Phi_z \to \Sigma \to D$).
A random phase diffuser creates depth-dependent PSFs through defocus-modulated wave propagation: $\text{PSF}(z) = |\mathcal{F}\{\exp(i\phi_{\text{diffuser}} + i \cdot \text{defocus}(z) \cdot r^2)\}|^2$.  Each depth plane produces a unique speckle pattern; the 2D sensor measurement is the sum of all convolved depth planes.  This is the $C$ primitive instantiated per depth plane.  PSNR: 20.3\,dB at 8:1 compression.

\begin{verbatim}
modality: 3d_lensless
carrier: photon
geometry: single_shot, 128x128, n_depths=8
object: 128x128x8 3D volume (sparse depth planes)
forward_model: Convolve_z(C, psf=diffuser(z)) -> Sum(Sigma) -> Detect(D)
noise: Poisson I_0=5000 + Gaussian sigma=3
target: PSNR >= 15 dB
system_elements:
  source: broadband LED
  optics: random phase diffuser (feature_scale=1.0, defocus_max=40)
  detector: CMOS 128x128
\end{verbatim}

\textbf{Temporal-coded lensless} ($M \to C \to \Sigma \to D$).
A binary DMD mask modulates each of 8 video frames before diffuser convolution and temporal integration during a single exposure.  The mask provides per-pixel temporal diversity.  PSNR: 31.6\,dB.

\textbf{Spectral lensless} ($M \to W \to C \to \Sigma \to D$).
A coded mask followed by prism dispersion separates 8 spectral bands before diffuser convolution and spectral integration.  PSNR: 36.3\,dB.

\subsection{S2.3. 4D Systems (16:1 compression)}

\textbf{4D Spectral-Depth} ($M \to W_\lambda \to \Phi_z \to \Sigma \to D$).
Combines coded mask, spectral dispersion, and diffuser depth encoding to recover a $(z, \lambda)$ datacube ($4 \times 4 = 16$ planes) from a single 2D shot.  PSNR: 23.9\,dB.

\textbf{4D Temporal DMD} ($M \to \Phi_z \to \Sigma \to D$).
DMD binary masks provide temporal coding; diffuser PSFs encode depth.  Recovers $(z, t)$ datacube at 16:1.  PSNR: 25.4\,dB.

\textbf{4D Temporal Streak} ($M \to W_t \to \Phi_z \to \Sigma \to D$).
Passive temporal dispersion (streak camera) replaces active DMD coding.  A fixed mask provides spatial diversity.  PSNR: 25.3\,dB.

\subsection{S2.4. 5D Systems (64:1 compression)}

\textbf{5D Full DMD} ($M \to W_\lambda \to \Phi_z \to \Sigma \to D$).
Full 5D recovery of $(x, y, z, \lambda, t)$ from a single shot at $4 \times 4 \times 4 = 64$:1 compression using coded mask + spectral dispersion + diffuser depth encoding.  PSNR: 29.9\,dB.

\textbf{5D Full Streak} ($M \to W_\lambda \to W_t \to \Phi_z \to \Sigma \to D$).
Replaces DMD with passive streak temporal dispersion.  The longest chain in the FPB framework (5 operators before detection).  PSNR: 29.9\,dB.

\subsection{S2.5. Design Patterns}

Three patterns emerge from the novel designs:

\begin{enumerate}[nosep]
\item \textbf{Active modulation (DMD) $>$ passive dispersion}: Binary masks provide per-pixel measurement diversity that continuous dispersion cannot match.  DMD-based temporal coding (31.6\,dB) outperforms streak-based (25.3\,dB) at equal compression.
\item \textbf{Algorithm choice is critical for compressive systems}: FISTA+TV or R-L dominate across novel modalities.  Inter-method PSNR CoV reaches 42--48\% for compressive systems vs.\ $<$6\% for well-conditioned ones.
\item \textbf{Graceful compression degradation}: 1:1 $\to$ 43.7\,dB; 8:1 $\to$ 20--36\,dB; 16:1 $\to$ 24--25\,dB; 64:1 $\to$ 29--30\,dB.
\end{enumerate}

% ========================================================================
\section{S3. Maskless Control Group}
% ========================================================================

To validate that spatial modulation ($M$) is a prerequisite for compressive dispersion-based imaging, we tested 9 configurations that combine depth encoding ($\Phi_z$) with spectral ($W_\lambda$) and/or temporal ($W_t$) dispersion but \emph{without} a coded mask.  All 9 configurations compile through the structural compiler (Gates~1--5) but are flagged by the Judge Agent (Gate~1, recoverability) as severely ill-conditioned:

\begin{itemize}[nosep]
\item Condition number $> 10^6$ for all maskless configurations
\item Reconstruction PSNR: 6--9\,dB (near noise floor)
\item Adding a binary coded mask ($M$) improves PSNR by 5--15\,dB at equal compression
\end{itemize}

This confirms a fundamental principle: dispersion operators ($W$) provide only global spectral/temporal diversity, while coded masks ($M$) provide per-pixel spatial diversity.  Both are necessary for well-conditioned compressive recovery.

% ========================================================================
\section{S4. Encoder Comparison for Depth Encoding}
% ========================================================================

Three physically realisable depth-encoding mechanisms were compared across 5 modulation combinations ($3 \times 5 = 15$ configurations total):

\textbf{Diffuser} ($\Phi_z^{\text{diff}}$): Random phase plate with defocus-dependent PSFs.  $\text{PSF}(z) = |\mathcal{F}\{\exp(i\phi_{\text{random}} + i \cdot \text{defocus}(z) \cdot r^2)\}|^2$.

\textbf{Multi-lens} ($L_z$): Microlens array producing depth-dependent parallax sub-pixel shifts.

\textbf{Meta-surface} ($\Psi_z$): Engineered phase pattern (cubic + trefoil + astigmatism aberrations) with defocus modulation.

\begin{table}[htbp]
\centering
\caption{Depth encoder comparison: best PSNR (dB) across 3 encoders $\times$ 5 modulation combinations.  $N_z{=}8$ depth planes, $N_\lambda{=}4$ spectral bands, $N_t{=}4$ temporal frames, image size $128{\times}128$.}
\small
\begin{tabular}{@{}lccccc@{}}
\toprule
\textbf{Encoder} & \textbf{3D only} & \textbf{+Mask} & \textbf{+M+$W_\lambda$} & \textbf{+M+$W_t$} & \textbf{+M+$W_\lambda$+$W_t$} \\
& (8:1) & (8:1) & (32:1) & (32:1) & (128:1) \\
\midrule
Diffuser & 17.0 & 20.3 & 24.6 & 26.9 & 38.1 \\
Multi-lens & 16.6 & \textbf{44.1} & \textbf{41.0} & \textbf{36.1} & \textbf{38.7} \\
Meta-surface & 16.6 & 16.2 & 24.6 & 25.8 & 38.0 \\
\bottomrule
\end{tabular}
\end{table}

\textbf{Key findings:}
\begin{itemize}[nosep]
\item Without coded mask: all three encoders perform similarly (16.6--17.0\,dB), limited by the ill-conditioning of depth-only recovery at 8:1 compression.
\item Multi-lens + mask dominates: parallax shifts are simple, well-conditioned operators (sub-pixel translations).  Combined with a binary mask's per-pixel diversity, the composite forward operator is nearly diagonal, yielding 44.1\,dB.
\item At maximal compression (128:1, 5D), all encoders converge to ${\sim}$38\,dB because the mask and dispersion operators dominate the measurement diversity.
\end{itemize}

\textbf{PSF diversity analysis} (mean inter-depth cross-correlation):
\begin{itemize}[nosep]
\item Diffuser: 0.29 (moderate diversity, structured speckle)
\item Multi-lens: 0.17 (lowest correlation, shift-based diversity)
\item Meta-surface: 0.46 (highest correlation, engineered but similar patterns)
\end{itemize}

% ========================================================================
\section{S5. Reconstruction Details}
% ========================================================================

\subsection{S5.1. Algorithms}

Five reconstruction algorithms were evaluated for each novel modality:

\textbf{Wiener deconvolution}: Frequency-domain least-squares with noise regularisation.  Fastest ($<$0.1\,s) but lowest quality for compressive systems.

\textbf{GAP-TV}: Generalised alternating projection with total variation regularisation~\cite{meng2020gap}.  Alternates between data projection and TV denoising.

\textbf{FISTA+TV}: Fast iterative shrinkage-thresholding with TV proximal operator~\cite{beck2009fista}.  Nesterov-accelerated gradient descent with $O(1/k^2)$ convergence.

\textbf{ADMM+TV}: Alternating direction method of multipliers with TV splitting~\cite{boyd2011admm}.  Variable splitting separates data fidelity from regularisation.

\textbf{Richardson-Lucy (R-L)}: Maximum-likelihood deconvolution for Poisson noise.  Multiplicative updates preserve non-negativity.

\subsection{S5.2. Forward Model Implementation}

All forward models are implemented as linear operators with explicit adjoint computation.  The measurement for a multi-dimensional modality with depth encoding, mask, and dispersion is:
\[
y = \sum_{z,\lambda,t} D\left[ \text{shift}_{(\Delta x_\lambda, \Delta y_t)} \left( M \odot \text{conv}(x_{z,\lambda,t}, \text{PSF}(z)) \right) \right] + n
\]
where $D$ is detection, $M$ is the binary mask, $\text{conv}$ denotes convolution with depth-dependent PSF, $\text{shift}$ implements spectral/temporal dispersion as lateral shifts, and $n$ is Poisson + Gaussian noise.

% ========================================================================
\section*{References}
% ========================================================================

\begingroup
\renewcommand{\section}[2]{}
\begin{thebibliography}{10}

\bibitem{paperII}
Yang, C.
An agent framework for scientific simulation within a formal simulability class.
\textit{Preprint} (2026).
Available at \url{https://github.com/integritynoble/Physics_World_Model}.

\bibitem{meng2020gap}
Meng, Z., Ma, J., \& Yuan, X.
End-to-end low cost compressive spectral imaging with spatial-spectral self-attention.
\textit{Proc.\ ECCV} (2020).

\bibitem{beck2009fista}
Beck, A. \& Teboulle, M.
A fast iterative shrinkage-thresholding algorithm for linear inverse problems.
\textit{SIAM J.\ Imaging Sci.} \textbf{2}, 183--202 (2009).

\bibitem{boyd2011admm}
Boyd, S., Parikh, N., Chu, E., Peleato, B. \& Eckstein, J.
Distributed optimization and statistical learning via the alternating direction method of multipliers.
\textit{Found.\ Trends Mach.\ Learn.} \textbf{3}, 1--122 (2011).

\end{thebibliography}
\endgroup

\end{document}
